\documentclass[10pt,a4paper]{amsart}
\usepackage[margin=19mm]{geometry}
\usepackage{amsmath,amssymb,mathtools}
\usepackage[scr=rsfs,cal=euler]{mathalfa}
\usepackage{xfrac}
\usepackage[unicode,hidelinks,bookmarks=false]{hyperref}
\newcommand{\ie}{i.e.\ }
\newcommand{\cf}{cf.\ }
\newcommand{\resp}{respectively\ }
\newcommand{\Subsection}{\subsection}
\providecommand{\nobreakdash}{\nobreak\hbox{-}}
\setlength{\emergencystretch}{2em}
\multlinegap0pt
\usepackage{bm}
\usepackage{bbm}
\usepackage{dsfont}
\usepackage{xspace}
\usepackage{cancel}
\usepackage{mathtools}
\usepackage[shortlabels]{enumitem}
\usepackage{amsthm}
\makeatletter
\@ifundefined{prop}{\newtheorem{prop}{Prop}}{}
\@ifundefined{rmk}{\newtheorem{rmk}{Rmk}}{}
\makeatother
\makeatletter
\newcommand{\CN}{E^{\it N}}
\newcommand{\Id}{{\rm Id}}
\newcommand{\bi}{\mathbf i}
\newcommand{\bj}{\mathbf j}
\newcommand{\bs}{\boldsymbol}
\newcommand{\cc}{{^\circ}}
\renewcommand{\d}{{\mathrm d}}
\newcommand{\defeq}{{\stackrel{\rm def}{=}}}
\newcommand{\la}{{\langle}}
\newcommand{\lmt}{\longmapsto}
\newcommand{\mc}{\mathcal}
\newcommand{\ra}{{\rangle}}
\expandafter\ifx\csname thebibliography\endcsname\relax
\renewenvironment{thebibliography}[1]
     {\section*{\refname}
      \@mkboth{\MakeUppercase\refname}{\MakeUppercase\refname}
      \begin{enumerate}[label={[\arabic{enumi}]},itemindent=*,leftmargin=2.5em]
      \@openbib@code
      \sloppy
      \clubpenalty4000
      \@clubpenalty \clubpenalty
      \widowpenalty4000
      \sfcode`\.\@m}
     {\def\@noitemerr
       {\@latex@warning{Empty `thebibliography' environment}}
      \end{enumerate}}
\fi
\newcommand{\two}{{\sqrt{2}}}
\makeatother
\title[Stochastic motions of the two-dimensional many-body delta-Bo]{Stochastic motions of the two-dimensional many-body delta-Bose gas, IV: Transformations of relative motions}
\author{Yu-Ting Chen}
\date{Source: arXiv:2401.17243v3 (CC BY 4.0)}
\begin{document}
\maketitle

\noindent\emph{Source and licence.} Stochastic motions of the two-dimensional many-body delta-Bose gas, IV: Transformations of relative motions. Version arXiv:2401.17243v3 (CC BY 4.0), distributed under the Creative Commons Attribution 4.0 International licence, as recorded in the arXiv licence metadata for this exact version and shown on the abstract page https://arxiv.org/abs/2401.17243v3. Full rights notice: licenses/arxiv-2401.17243-CC-BY-4.0.txt.\par\smallskip

\noindent\textit{Editorial scope.} Counted unit: the Rmk whose source label is \texttt{rmk:com} in the article (source0.tex lines 515--524, source label \texttt{rmk:com}), delivered together with the article's own complete original proof of it (source0.tex lines 526--601, one \texttt{proof} environment of 76 lines). No mathematics is written, repaired or re-derived: statement and proof are the article's own text line for line. Required context retained uncounted: prop:RM (lines 430--507); def:R (lines 448--453); def:Rinv (lines 456--463); def:Zi (lines 468--470); def:tZ (lines 479--481); def:bijective (lines 495--504). Every definition, display and label of those lines is retained unchanged. Omitted from this package and not redistributed: 1--429, 454--455, 464--467, 471--478, 482--494, 505--514, 525--525, 602--769. Nothing inside a retained excerpt is omitted or altered: every retained body is delivered exactly as the article's lines are written, so no omission receipt of any kind accompanies this package. Citations: the retained excerpts carry no citation command at all, so no bibliography entry of the article is transcribed and no reference is redistributed. Reference adaptation: none was needed and none was made; every reference inside the delivered unit resolves inside it, so no unresolved reference or citation is produced. Printed slips kept literal: no printed word, symbol, hypothesis or number of the counted statement or of its proof was altered. Layout and class: the article's own class is \texttt{article} with packages including enumitem, tocloft, titletoc, pdfpages, amsmath, amssymb, amsthm, leftidx, epsfig, graphics, graphicx, colordvi, mathrsfs, stmaryrd; the wrapper is the lane's pinned \texttt{amsart} editorial wrapper at 10pt on A4, which restates the theorem-like environments the retained text needs and adds the article's own macro definitions that the retained text needs (\texttt{\textbackslash CN}, \texttt{\textbackslash Id}, \texttt{\textbackslash bi}, \texttt{\textbackslash bj}, \texttt{\textbackslash bs}, \texttt{\textbackslash cc}, \texttt{\textbackslash d}, \texttt{\textbackslash defeq}, \texttt{\textbackslash la}, \texttt{\textbackslash lmt}, \texttt{\textbackslash mc}, \texttt{\textbackslash ra}, \texttt{\textbackslash two}), with extra wrapper packages (bm, bbm, dsfont, xspace, cancel, mathtools, enumitem). The delivered numbering is the unit's own. Assets: no retained span contains a figure environment, a graphics-inclusion command, a PDF-page inclusion, a table or a TikZ picture; no image file is copied into this package, the package declares no asset dependency and no image file is redistributed with it. Licence: version arXiv:2401.17243v3 of this article is distributed under the Creative Commons Attribution 4.0 International licence, as recorded in the arXiv licence metadata for this exact version and shown on the abstract page https://arxiv.org/abs/2401.17243v3. A full rights notice is delivered at licenses/arxiv-2401.17243-CC-BY-4.0.txt. Characters: every retained excerpt is pure ASCII, so the delivered engine, XeLaTeX with its default Latin Modern text font, reports no missing character.
\begin{prop}[Main result]\label{prop:RM}
{\rm (1$\cc$)} As vector spaces,
$\CN$ is isomorphic to  $\delta \CN\times E$ via the following
linear transformation $\mathsf R:\CN\to \delta \CN\times E$:
\begin{align}
\mathsf R\;\defeq\, \begin{bmatrix}
-1/\two & 1/\two & 0&\cdots &0\\
0 &-1/\two&1/\two& \cdots &0\\
\vdots &\vdots &\ddots &\vdots&\vdots\\
0 &0 &\cdots &-1/\two &1/\two\\
1/N&1/N &\cdots &1/N &1/N
\end{bmatrix}_{N\times N}:
\begin{bmatrix}
z^1\\
z^2\\
\vdots\\
z^N
\end{bmatrix}
\lmt \begin{bmatrix}
z^{(2,1)}\\
\vdots\\
z^{(N,N-1)}\\
z^\Sigma
\end{bmatrix}.\label{def:R}
\end{align}
Specifically, $\mathsf R$ is invertible, and the matrix inverse admits the following formula:
\begin{align}\label{def:Rinv}
(\mathsf R^{-1})_{ij}=
\begin{cases}
\displaystyle -\two (N-j)/N,&\; i\leq j,\; j\leq N-1,\\
\displaystyle \two j/N,&\;i>j,\;j\leq N-1,\\
\displaystyle 1, &\;1\leq i\leq N,\;j=N.
\end{cases}
\end{align} 

\noindent {\rm (2$\cc$)} There is a pathwise one-to-one correspondence between the two classes [$N$] and [$\mathcal E_N$] of families of processes $\{\mc Z^i_t \}_{1\leq i\leq N}$ and $\{\mathcal Z^\bi_t \}_{\bi\in\mc E_N}\cup \{\mc Z^\Sigma_0+W^\Sigma_t \}$ defined as follows.
\begin{itemize}[leftmargin=1.2cm]
\item [{\bf [$\bs N$]}] The processes $\{\mc Z^i_t \}$  satisfy the following SDEs:
\begin{align}\label{def:Zi}
\mc Z^i_t=\mc Z^i_0-\sum_{j:j\neq i}\int_0^t a_{i j}(s)(\mc Z^i_s-\mc Z^j_s)\d s+W^i_t,\quad \forall\;1\leq i\leq N,
\end{align}
such that the following conditions hold:
\begin{itemize}
\item [$\bullet$] For all $i\neq j$, $a_{ij}(s)=a_{ij}(s,\omega)$ is progressively measurable, the symmetry $a_{ij}=a_{ji}$ holds, and
$\int_0^t |a_{ij}(s)(\mc Z^i_s-\mc Z^j_s)|\d s<\infty$.
\item [$\bullet$]  $\{W^i_t \}_{1\leq i\leq N}$ consists of independent $d$-dimensional standard Brownian motions starting from $0$.
\end{itemize}

\item [{\bf [$ \mathbfcal{E}_{\bs N}$]}] The processes $\{\mc Z^\bi_t \}$ satisfy the following SDEs:
\begin{align}\label{def:tZ}
\mc Z^\bi_t=\mc Z^\bi_0-\sum_{\bj\in \mc E_N}\sigma(\bi)\cdot \sigma(\bj)\int_0^t a_\bj(s)\mc Z^\bj_s\d s+W^{\bi}_t,\quad \forall\;\bi\in \mc E_N,
\end{align}
$\{\sqrt{N}W^\Sigma_t \}$ is a $d$-dimensional standard Brownian motion starting from $0$, and $\mc Z^\Sigma_0\in E$, such that the following conditions hold:
\begin{itemize}
\item [$\bullet$] $\{\mc Z^\bi_0\}_{\bi\in \mc E_N}$ is difference-consistent.
\item [$\bullet$] $\sigma(\bi)\in \{-1,0,1\}^N$ is such that the $i\prime$-th component is $1$, the $i$-th component is $-1$, and all the other components are zero. Hence, $\sigma(\bi)\cdot\sigma(\bj)$ in \eqref{def:tZ} denotes the usual dot product of vectors.
\item [$\bullet$]  For all $\bi\in \mc E_N$, $a_{\bi}(s)=a_{\bi}(s,\omega)$ are progressively measurable, and we have $\int_0^t |a_{\bi}(s)\mc Z^\bi_s|\d s<\infty$.
\item [$\bullet$] $\{W^\bi_t \}_{\bi\in\mathcal E_N}$ satisfies the following properties. 
(i) It is independent of $\{W^\Sigma_t \}$; (ii) it has the same law of a linear transformation of Brownian motion such that $\{W^\bi_t \}$ is a $d$-dimensional standard Brownian motion with zero initial condition $0$ and real-valued components $\{W_t^{\bi,\ell}\}$, $1\leq \ell\leq d$, such that the following identities of quadratic variation hold:
$\la W^{\bi,\ell},W^{\bj,\ell}\ra_t=\sigma(\bi)\cdot \sigma(\bj)t/2$ and $
\la W^{\bi,k},W^{\bj,\ell}\ra_t=0$
for all $1\leq k, \ell\leq d$ with $k\neq \ell$ and all $\bi,\bj\in \mc E_N$.
\end{itemize}
\end{itemize}
Specifically, the one-to-one correspondence between class [$N$] and class [$\mathcal E_N$] is defined as follows: 
\begin{align}\label{def:bijective}
\begin{cases}
\displaystyle \frac{\mc Z_t^{i\prime}-\mc Z_t^i}{\two}=\mathcal Z_t^\bi,
\quad a_{ij}(s,\omega)=a_{i\vee\!\wedge j}(s,\omega),\quad \frac{W^{i\prime}_t-W^i_t}{\two}=W^{\bi}_t,\\
\vspace{-.4cm}\\
\displaystyle
\frac{1}{N}\sum_{i=1}^N \mc Z^i_t=\mc Z^\Sigma_0+W^\Sigma_t,\quad 
\frac{1}{N}\sum_{i=1}^N W^i_t=W^\Sigma_t,
\end{cases}
\end{align}
where $i\!\vee\!\!\wedge j=j\!\vee\!\!\wedge i\,\defeq\, 
(\max\{i,j\},\min\{i,j\})\in \mc E_N$ for any integers $1\leq i\neq j\leq N$.
\end{prop}

\begin{rmk}[Symmetry and stochastic center of mass]\label{rmk:com}
{\rm (1$\cc$)} The mapping between $a_{i\vee\!\wedge j}$ and $a_{ij}$ in \eqref{def:bijective} is well-defined because of the symmetry assumption $a_{ij}=a_{ji}$.\medskip 

\noindent {\rm (2$\cc$)} By $a_{ij}=a_{ji}$, \eqref{def:Zi} implies that the {\bf stochastic center of mass} $N^{-1}\sum_{i=1}^N \mc Z^i_t$ satisfies
\begin{align}\label{eq:com}
\frac{1}{N}\sum_{i=1}^N \mc Z^i_t=\frac{1}{N}\sum_{i=1}^N \mc Z^i_0+\frac{1}{N}\sum_{i=1}^NW^i_t,
\end{align}
and $\{\sqrt{N}\cdot N^{-1}\sum_{i=1}^NW^i_t \}$ is a $d$-dimensional standard Brownian motion independent of the correlated Brownian motions $\{(W^{i\prime}_t-W^i_t)/\two \}_{\bi\in \mc E_N}$. In particular, it is according to \eqref{eq:com} that we choose the last row of $\mathsf  R$ in \eqref{def:R}. 
\qed 
\end{rmk}

\begin{proof}[Proof of Proposition~\ref{prop:RM} (1$\cc$)]
To prove \eqref{def:Rinv}, we make two observations to simplify the algebra a bit.
First, observe that by \eqref{def:R},
\begin{align}\label{def:Q}
\mathsf R=\begin{bmatrix}
1/\two &0&\cdots &0\\
0 &1/\two&\cdots &0\\
\vdots &\vdots &\ddots &\vdots\\
0 &\cdots & 1/\two &0\\
0 &\cdots &0 &1/N
\end{bmatrix} 
\mathsf  Q,\quad \mbox{for }\mathsf Q\;\defeq\, \begin{bmatrix}
-1 & 1 & 0&0&\cdots &0\\
0 &-1&1& 0&\cdots &0\\
\vdots &\vdots &\vdots &\ddots &\vdots&\vdots\\
0 &0 &0&\cdots &-1 &1\\
1&1 &1&\cdots &1 &1
\end{bmatrix}.
\end{align}
Second, for any $N\times N$-matrix $\widetilde{\mathsf M}=[\widetilde{\mathsf M}_1,\cdots,\widetilde{\mathsf M}_N]$ and any scalars $D_1,\cdots,D_N$,
\begin{align}\label{MD}
[\widetilde{\mathsf M}_1,\cdots,\widetilde{\mathsf M}_N]\mathsf {\rm diag}(\mathsf D_1,\cdots,\mathsf D_N)=[\mathsf D_1\widetilde{\mathsf M}_1,\cdots,\mathsf D_N \widetilde{\mathsf M}_N].
\end{align}
Thus, the proposed inverse in \eqref{def:Rinv} takes the form of the left-hand side of \eqref{MD} with $\mathsf D_1=\cdots=\mathsf D_{N-1}=\sqrt{2}$, $\mathsf D_N=N$, and $\widetilde{\mathsf M}$ chosen to be $\mathsf M$ given by 
\begin{align}\label{def:Qinv}
\mathsf M_{ij}\defeq
\begin{cases}
\displaystyle -(N-j)/N,&\; i\leq j,\; j\leq N-1,\\
\displaystyle j/N,&\;i>j,\;j\leq N-1,\\
\displaystyle 1/N, &\;1\leq i\leq N,\;j=N.
\end{cases}
\end{align} 
Combining these two observations shows that it is enough to prove the invertibility of $\mathsf Q$ with $\mathsf M=\mathsf Q^{-1}$, or just the identity
$\mathsf Q\mathsf M=\Id $. 

Now, to verify $\mathsf Q\mathsf M=\Id $, we use the definitions of $\mathsf Q$ and $\mathsf M$
 in \eqref{def:Q} and \eqref{def:Qinv} and consider $(\mathsf Q\mathsf M)_{ij}=\sum_{k=1}^N \mathsf Q_{ik}\mathsf M_{kj}$, $1\leq i,j\leq N$, for five cases which are mutually exclusive, collectively exhaustive: \medskip 

\noindent {\bf (1)} For all $1\leq i< j\leq N$,
\begin{align*}
(\mathsf Q\mathsf M)_{ij}&
=\mathsf Q_{ii}\mathsf M_{ij}+\mathsf Q_{i(i+1)}\mathsf M_{(i+1)j}
=-\mathsf M_{ij}+\mathsf M_{(i+1)j}
=
\begin{cases}
\displaystyle \frac{N-j}{N}-\frac{N-j}{N}=0,&j\leq N-1,\\
\vspace{-.3cm}\\
\displaystyle -\frac{1}{N}+\frac{1}{N}=0,&j=N.
\end{cases}
\end{align*}

\noindent {\bf (2)} For all $1\leq j<i\leq N-1$,
\begin{align*}
(\mathsf Q\mathsf M)_{ij}&
=\mathsf Q_{ii}\mathsf M_{ij}+\mathsf Q_{i(i+1)}\mathsf M_{(i+1)j}=-\mathsf M_{ij}+\mathsf M_{(i+1)j}=-\frac{j}{N}+\frac{j}{N}=0.
\end{align*}

\noindent {\bf (3)} For $1\leq j<i=N$,
\begin{align*}
(\mathsf Q\mathsf M)_{Nj}
=\sum_{k=1}^N\mathsf M_{kj}=j \left(-\frac{N-j}{N}\right)+(N-j)\frac{j}{N}=0.
\end{align*}

\noindent {\bf (4)} For all $1\leq i=j\leq N-1$,
\begin{align*}
(\mathsf Q\mathsf M)_{ii}
=\mathsf Q_{ii}\mathsf M_{ii}+\mathsf Q_{i(i+1)}\mathsf M_{(i+1)i}=-\mathsf M_{ii}+\mathsf M_{(i+1)i}=\frac{N-i}{N}+\frac{i}{N}=1.
\end{align*}

\noindent (5) For $i=j=N$,
\begin{align*}
(\mathsf Q\mathsf M)_{NN}
=\sum_{k=1}^N\mathsf M_{kN}=\sum_{k=1}^N \frac{1}{N}=1.
\end{align*}
By the last equalities in these five cases,  we have proved that the matrix $\mathsf M$ defined by \eqref{def:Qinv} satisfies $\mathsf Q\mathsf M={\rm Id}$. The proof is complete.
\end{proof}

\end{document}
